% == == == == == == == == == == == == == == == == == == == == == == == == == ==
    == == == == == == == == == == == ==
    = % Section 6 : Work Breakdown Structure(WBS) % == == == == == == == == ==
      == == == == == == == == == == == == == == == == == == == == == == == == ==
      == == == ==
    =

\section {
  Work Breakdown Structure
}
\label {
sec:
  wbs
}
\sectionauthors { CM, MN }

\subsection{WBS Overview}

The Work Breakdown
Structure(WBS) breaks the project scope into manageable work packages \parencite{wbs_guide}. It organizes all work under six main categories: Research, Design, Build, Test, Close, and Project Management Overhead. This structure adds detail at every level to clearly define specific tasks and deliverables.

\subsubsection{Team Member Assignments}

Each leaf-level activity is assigned a primary owner who is accountable for its completion.

\subsubsection{Ownership Structure}

Activity ownership follows a RACI-based(Responsible Accountable Consulted Informed) assignment model \parencite{pmi_raci}:

\begin{
  itemize
}
    \item \textbf{Owner (A/R)}: Primary team member accountable for and responsible for completing the activity. Shown in the ``Owner'' column of WBS tables.
    \item \textbf{Sub-Owners}: Supporting team members who contribute as Responsible (R), Consulted (C), or Informed (I) parties.
\end{
      itemize
    }

    \noindent Ownership ensures clear accountability,
        balanced workload distribution, and skills alignment. \textbf {
      For detailed resource allocation including
              secondary assignments and person -
          day calculations,
          see the WBS Activity Assignment List in Section \
      ref{sec : resource_planning},
          starting in Table \ref{tab : activity_assignments}.
    }

    \subsection{Control Accounts}

    The WBS is organized under six primary control accounts \parencite{
        dau_control_account} :

\begin{table}[H]
\centering
\caption {
      WBS Control Account Summary
    }
    \label {
    tab:
      wbs_control_accounts
    }
    \begin{tabular} {
      @{} llrl @{}
    }
    \toprule
\textbf{ID} & \textbf{Control Account} & \textbf{Tasks} & \textbf {
      Description
    }
    \\
\midrule 1.1 & Research & 13 & Requirements gathering,
        technology evaluation \ 1.2 & Design & 24 & Architecture, mechanical,
        electrical, software design \ 1.3 & Build & 20 & Procurement,
        fabrication, assembly, integration \ 1.4 & Test & 15 & Validation,
        verification, system testing \ 1.5 & Close & 7 & Documentation, handoff,
        final report \ 1.6 & PM Overhead & 5 & Meetings, status reports,
        risk management \\
\midrule & \textbf{Total} & \textbf{84} & \\
\bottomrule
\end {
      tabular
    }
    \end { table }

    \subsection{WBS Hierarchy}

    The hierarchical structure follows this pattern :

\begin {
      itemize
    }
    \item \textbf{Level 1} : Project Root(1.0 Project Name)
    \item \textbf{Level 2}
        : Control Accounts(Research, Design, Build, Test, Close, PM Overhead)
    \item \textbf{Level 3}
        : Categories(Requirements Analysis, System Architecture)
    \item \textbf{Level 4}
        : Work Packages(Functional Requirements, Algorithm Design)
    \item \textbf{Level 5} : Tasks(Document user stories, Evaluate options)
\end {
      itemize
    }

    \subsection{WBS Diagram}

    The complete WBS overview diagram is included in Appendix \
    ref{app : wbs}(Figure \ref{app : wbs})
        .The following
        subsections show the detailed WBS structure in table format,
        generated from the project's activities database.

                % Native LaTeX WBS tables(auto - generated from activities.yaml)
\input{wbs / wbs_complete}

                % Phase
            - specific WBS breakdowns with Forest diagrams and tables
\input {
      wbs / wbs_research
    }
    \input { wbs / wbs_design }
    \input { wbs / wbs_build }
    \input { wbs / wbs_test }
    \input { wbs / wbs_close }

    \subsection{Activity Naming Convention}

        All WBS activities follow the noun -
        verb format as required :

\begin {
      itemize
    }
    \item \textbf{Correct} : ``Research algorithms'', ``Design chassis'', ``Test system''
    \item \textbf{Incorrect} : ``Algorithm research'', ``Chassis design'', ``Testing''
\end {
      itemize
    }

    \subsection{Duration Guidelines}

    Per FTP requirements,
        all leaf - level activities are limited to a maximum of 5 days
                       .Activities originally estimated at longer durations have
                           been decomposed into smaller tasks.
